\documentclass[runningheads]{llncs}

\usepackage{amsmath,amssymb}
\usepackage{booktabs}
\usepackage{multirow}
\usepackage{graphicx}
\usepackage{xcolor}
\usepackage{url}
\usepackage{hyperref}

\newcommand{\tbd}{\textcolor{red!70!black}{\textbf{TBD}}}
\newcommand{\CUG}{\mathrm{CUG}}
\newcommand{\PVA}{\mathrm{PVA}}
\newcommand{\GUD}{\mathrm{GUD}}
\newcommand{\CEG}{\mathrm{CEG}}
\newcommand{\IOD}{\mathrm{IOD}}

% Result figures compile as placeholders until generated by the analysis pipeline.
\newcommand{\resultfigure}[3]{%
  \begin{figure}[t]
    \centering
    \IfFileExists{#1}{\includegraphics[width=#2\linewidth]{#1}}{%
      \fbox{\parbox[c][34mm][c]{0.92\linewidth}{\centering
      \textbf{Result figure pending experiment execution}\\[2mm]
      Expected artifact: \texttt{#1}}}}
    \caption{#3}
  \end{figure}
}

\begin{document}

\title{Separating Personalization from Discrimination in Personality-Aware LLM Recommendations}
\titlerunning{Personalization vs. Discrimination in LLM Recommendation}

% ECIR 2027 full papers are double blind.
\author{Anonymous Authors}
\authorrunning{Anonymous Authors}
\institute{Anonymous Institution}

\maketitle

\begin{abstract}
Large language models (LLMs) increasingly act as conversational recommenders, yet auditing their fairness is difficult because a recommendation that changes after receiving additional user context is not necessarily unfair: the change may be legitimate personalization. Existing LLM-recommendation audits often conflate ranking sensitivity with discrimination, use weakly grounded personality descriptions, or evaluate too few models and domains to establish reliable conclusions. We introduce a preference-conditioned evaluation design that separates \emph{behavioral change} from its \emph{consequence for user utility and exposure}. The benchmark combines real interaction histories with controlled demographic counterfactuals and measured Big Five personality profiles, including shuffled-personality and one-trait counterfactual controls. We plan a six-dataset, six-LLM-family evaluation with candidate-constrained ranking, repeated generations, explicit invalid-output accounting, and utility--fairness analysis. We further evaluate lightweight prompting and preference-aligned counterfactual re-ranking to reduce harmful identity-conditioned gaps without erasing useful personalization. Numerical results are intentionally omitted from this development draft and will be generated from the frozen experiment pipeline before submission.
\keywords{LLM recommendation \and fairness \and personalization \and personality \and counterfactual evaluation}
\end{abstract}

\section{Introduction}
Large language models have created a conversational recommendation paradigm in which users can express preferences, constraints, and personal context in natural language. This flexibility also creates a difficult fairness problem. A recommendation list may change after the model receives a user's age, gender, country, or personality information; however, \emph{change alone does not identify harm}. Some changes may improve preference alignment, while others may encode stereotypes or create unequal recommendation quality.

Early RecLLM fairness benchmarks measured recommendation sensitivity to protected attributes and helped expose important risks \cite{zhang2023fairllm}. Subsequent work has shown why sensitivity is not a sufficient fairness criterion. CFaiRLLM explicitly argues that list discrepancies can arise from genuine preference-aligned personalization \cite{deldjoo2025cfairllm}; work on chatbot recommendations similarly finds that identity can influence recommendations in ways that mix personalization and racial stereotyping \cite{kantharuban2025stereotype}. LLMs may also infer demographic traits implicitly from conversational cues, causing stereotype-driven personalization even when attributes are not stated \cite{neplenbroek2025reading}.

This paper starts from a stricter question: \emph{holding preference evidence, candidate items, task, and prompt syntax fixed, what recommendation effects are caused by demographic or personality context, and are those effects beneficial, neutral, or harmful with respect to held-out user preferences?} We distinguish demographics from psychographics rather than treating personality as another sensitive demographic attribute. Personality experiments use measured Big Five/OCEAN information and matched negative controls, while demographic experiments use paired counterfactual interventions.

We ask four research questions:

\begin{description}
\item[RQ1:] \textbf{Personalization vs. discrimination.} When identity or personality context changes an LLM ranking, is the change preference-aligned, or does it create an unjustified counterfactual utility/exposure gap?
\item[RQ2:] \textbf{Grounded personality.} Do measured Big Five traits provide recommendation value beyond preference history, and do such gains survive shuffled-personality and one-trait counterfactual controls?
\item[RQ3:] \textbf{Generalization and reliability.} How stable are the conclusions across six datasets, six LLM families, prompt templates, ranking cutoffs, and repeated generations?
\item[RQ4:] \textbf{Mitigation.} Can inexpensive prompting and preference-aligned counterfactual re-ranking reduce harmful gaps while preserving utility and useful personality-driven personalization?
\end{description}

Our contributions are therefore a \emph{problem reformulation}, a \emph{grounded benchmark design}, a \emph{reliability protocol}, and a \emph{utility-aware mitigation evaluation}. We do not claim that any model is fairer, or that personality improves recommendation, until the frozen experiments produce evidence for those claims.

\section{Related Work}
\subsection{Fairness in LLM-Based Recommendation}
FaiRLLM established an early benchmark for evaluating sensitive-attribute effects in ChatGPT recommendations \cite{zhang2023fairllm}. UP5 investigates user-side fairness in foundation-model recommendation and proposes counterfactually fair prompting \cite{hua2024up5}. CFaiRLLM moves evaluation toward consumer utility by warning that ranking differences need not be bias when they reflect true preferences \cite{deldjoo2025cfairllm}. Our work takes this distinction as a starting point rather than a novelty claim: we add measured personality, matched negative controls, candidate-constrained held-out utility, multi-source reliability analysis, and mitigation that explicitly attempts to retain beneficial personalization.

\subsection{Personalization, Identity, and Stereotypes}
Identity-conditioned recommendation is especially difficult because personalization and stereotyping can produce superficially similar output changes. Kantharuban et al. demonstrate that user identity significantly affects chatbot recommendations and that racial stereotypes can appear under both explicit and implicit identity cues \cite{kantharuban2025stereotype}. Neplenbroek et al. further show that LLMs infer demographic information from subtle cues and that such implicit personalization can lower response quality for minority users \cite{neplenbroek2025reading}. These findings motivate our requirement that demographic counterfactuals be evaluated against stable preference evidence rather than through list overlap alone.

\subsection{Personality-Aware Recommendation}
Personality-aware recommendation predates LLMs. The GroupLens Personality 2018 resource links recommender behavior and satisfaction to measured personality \cite{nguyen2018personality}, while music recommendation work has collected BFI/BFI-2 profiles and item ratings \cite{klec2023bigfive}. We use such psychometrically grounded sources to test whether personality contributes predictive value rather than inserting informal descriptors such as ``introverted engineer'' into synthetic prompts.

\section{Problem Formulation}
Let user $u$ have observed preference history $H_u$, auditable candidate set $C_u$, and held-out relevance labels $Y_u \subseteq C_u$. Let $z_u$ denote optional contextual information. An LLM recommender $M$ is stochastic and induces a distribution over rankings,
\begin{equation}
R_{u}^{z,s} \sim p_M(R \mid H_u, C_u, z, t, s),
\end{equation}
where $t$ is a fixed prompt template and $s$ indexes repeated generations or provider-supported randomization.

Let $U(R,Y)$ be a conventional recommendation-utility function such as nDCG@K or Recall@K. Let $D(R,R')$ measure behavioral ranking change, such as $1-\mathrm{RBO}$ or $1-\mathrm{Jaccard}$. Crucially, $D$ is \emph{not} a fairness metric by itself.

For a paired demographic intervention $a \rightarrow a'$ that leaves $H_u$, $C_u$, and $t$ unchanged, we define the signed per-user Counterfactual Utility Gap,
\begin{equation}
\CUG_u(a,a') = U(R_u^{a},Y_u)-U(R_u^{a'},Y_u).
\end{equation}
We report the paired distribution, absolute gap, confidence interval, and effect size rather than reducing the experiment to a single list-similarity score.

When reliable item-group metadata is available, we analogously compare discounted exposure distributions and compute a Counterfactual Exposure Gap $\CEG$. Group Utility Disparity $\GUD$ summarizes differences in mean utility across pre-specified groups, while Invalid Output Disparity $\IOD$ treats persistent generation failures as observable system behavior rather than silently discarding them.

For measured personality, let $p_u$ be the true OCEAN vector and $\pi(p)_u$ a shuffled profile taken from another user under a fixed randomization scheme. The Personality Value Added statistic is
\begin{equation}
\PVA_u = U(R_u^{p_u},Y_u)-U(R_u^{\pi(p)_u},Y_u).
\end{equation}
A positive aggregate $\PVA$ is evidence that the measured profile contains user-specific recommendation value beyond merely adding personality-shaped text. We additionally intervene on one OCEAN dimension at a time while holding the remaining profile fixed.

\section{Benchmark and Evaluation Design}
\subsection{Datasets}
We use two complementary tracks. The personality-grounded track contains datasets with measured psychometric information; the generalization track broadens domains and supports controlled counterfactual stress tests.

\begin{table}[t]
\centering
\caption{Planned six-dataset benchmark. Final post-filter sample counts and hashes will be generated by dataset adapters. ``Observed'' means supplied by the dataset, not inferred by us.}
\label{tab:datasets}
\small
\begin{tabular}{p{2.6cm}p{1.45cm}p{3.2cm}p{3.1cm}}
\toprule
Dataset & Domain & Grounding & Primary role \\
\midrule
Personality 2018 & Movies & interactions + measured Big Five & personality-grounded utility \\
Music Master / BFI-2 & Music & ratings + BFI/BFI-2 & personality/facet controls \\
REASONER & Short video & interactions + CBF-PI-15 responses & personality generalization \\
MovieLens-1M & Movies & interactions + dataset demographics & demographic counterfactuals \\
Last.fm & Music & listening histories; profile fields in selected release & demographic/domain generalization \\
MIND \cite{wu2020mind} & News & real click histories + logged impressions & candidate-constrained stress test \\
\bottomrule
\end{tabular}
\end{table}

MIND contains anonymized behavior rather than demographic identity; therefore any identity manipulation on MIND is explicitly treated as a synthetic counterfactual stress test, not an observed-demographic analysis. We never infer protected attributes from names, text, or embeddings.

\subsection{Candidate-Constrained Tasks}
Open-ended requests for ``25 movies'' confound fairness with model memorization, hallucination, item availability, and popularity. Each benchmark instance instead contains a preference history, an auditable candidate set, and held-out relevant items. The model must return exactly $K$ unique candidate IDs in structured JSON. The primary cutoff is $K=10$, with $K\in\{5,10,20\}$ as a sensitivity analysis.

\subsection{Controlled Context Conditions}
We use a fixed field-based prompt in which absent context is explicitly represented as \texttt{unspecified}. Conditions include: preference-only ($C_0$), observed demographics when available ($C_1$), paired demographic counterfactuals ($C_2$), true measured personality ($C_3$), shuffled personality ($C_4$), one-trait personality counterfactuals ($C_5$), and a pre-specified intersectional condition ($C_6$) on datasets that support it. Candidate order is randomized once per benchmark instance and then frozen across all paired conditions.

\subsection{Model Panel and Stochasticity}
We evaluate one pinned model snapshot from each of six independent families: OpenAI, Anthropic Claude, Google Gemini, DeepSeek, Alibaba Qwen, and Meta Llama. Exact identifiers, provider revisions, request timestamps, decoding parameters, prompt hashes, and code commit hashes are recorded in a machine-readable manifest. The main evaluation uses three generations per prompt; a stratified 10\% subset receives ten generations for a focused stochasticity audit. If a provider does not expose deterministic seed control, we report that limitation rather than assuming cross-provider seed equivalence.

\subsection{Metrics and Statistical Analysis}
Recommendation effectiveness is measured primarily by nDCG@10 and Recall@10. RBO and Jaccard quantify behavioral change only. Fairness analyses report paired CUG, group utility disparity, counterfactual exposure gaps when item metadata permits, and invalid-output disparities. Personality analysis reports $\PVA$ and trait-stratified effects. We pre-specify paired bootstrap confidence intervals, paired nonparametric/permutation tests, Holm correction within RQ families, effect sizes, and mixed-effects or variance-component analysis when supported by the final design.

\subsection{Output Validation}
Every response is parsed against the same schema, checked for exact $K$, uniqueness, and candidate membership. At most one format-only repair is allowed, which may serialize existing choices but may not add, remove, replace, or reorder them. Persistent failure remains in the dataset and contributes to $\IOD$.

\section{Mitigation: Preference-Aligned Identity Re-ranking}
RQ4 compares the unmitigated model with (i) an instruction that task-irrelevant demographic identity must not serve as a preference proxy, (ii) a counterfactually fair prompting baseline motivated by prior work \cite{hua2024up5}, and (iii) a lightweight re-ranker, \textbf{PAIR} (Preference-Aligned Identity Re-ranking).

Given neutral and counterfactual rankings, PAIR rewards candidates that receive consistent preference-based support while penalizing identity-specific exposure instability:
\begin{equation}
\mathrm{score}(i)=\mathrm{RelConsensus}(i)-\lambda\,\mathrm{CFInstability}(i).
\end{equation}
The trade-off parameter $\lambda$ is selected using validation instances only. We report a utility--fairness Pareto frontier rather than choosing a post-hoc operating point that makes one method look best.

\section{Results}
This development manuscript contains no invented numerical findings. The final section will be generated from frozen experiment outputs.

\subsection{RQ1: When Is a Ranking Change Personalization Rather Than Harm?}
\begin{table}[t]
\centering
\caption{RQ1 primary result structure. Values will be generated from paired user instances.}
\label{tab:rq1}
\small
\begin{tabular}{lcccc}
\toprule
Model family & $\Delta$ nDCG@10 & $|\CUG|$ & $\CEG$ & RBO@10 \\
\midrule
OpenAI & \tbd & \tbd & \tbd & \tbd \\
Anthropic & \tbd & \tbd & \tbd & \tbd \\
Google & \tbd & \tbd & \tbd & \tbd \\
DeepSeek & \tbd & \tbd & \tbd & \tbd \\
Qwen & \tbd & \tbd & \tbd & \tbd \\
Llama & \tbd & \tbd & \tbd & \tbd \\
\bottomrule
\end{tabular}
\end{table}

\resultfigure{figures/rq1_quadrant.pdf}{0.93}{RQ1 planned diagnostic: paired cases positioned by ranking change and held-out utility consequence. Ranking sensitivity with improved utility is separated from identity-conditioned quality loss.}

\subsection{RQ2: Does Measured Personality Add User-Specific Value?}
\resultfigure{figures/rq2_personality_forest.pdf}{0.95}{RQ2 planned forest plot of true-versus-shuffled personality effects ($\PVA$), stratified by dataset/model and accompanied by uncertainty intervals.}

\subsection{RQ3: Are Conclusions Reliable Across Models, Domains, and Generations?}
\resultfigure{figures/rq3_variance.pdf}{0.94}{RQ3 planned reliability analysis separating model-family, dataset, prompt-template, and repeated-generation variation.}

\subsection{RQ4: Can We Mitigate Harm Without Erasing Personalization?}
\resultfigure{figures/rq4_pareto.pdf}{0.94}{RQ4 planned utility--fairness Pareto frontier comparing unmitigated ranking, prompting baselines, and PAIR.}

\section{Discussion, Limitations, and Ethical Scope}
Our benchmark is designed to identify controlled recommendation disparities; it does not by itself establish unlawful discrimination or user-perceived harm. Counterfactual identity edits are interventions in a benchmark prompt, not claims that identity is mutable in the social world. Measured personality introduces separate ethical concerns: psychographic information can improve personalization but may also enable manipulation or intrusive profiling. Accordingly, the study evaluates whether personality provides incremental value and does not argue that production systems should collect personality by default.

The six datasets vary in age, sparsity, domains, licenses, and available user metadata. Cross-dataset aggregation must therefore preserve dataset-level effects. API-hosted LLMs may change over time, making exact version logging essential. Candidate-constrained ranking improves evaluability but differs from unconstrained conversational recommendation, so conclusions should not automatically transfer to open-ended generation.

For an IR-for-Good framing, the intended positive outcome is not simply ``more fairness research.'' The intervention path is: detect preference-irrelevant identity effects with an auditable benchmark; distinguish them from helpful personalization; expose model/domain-specific failure modes; and test mitigations that reduce harmful disparities without sacrificing relevance. This mechanism should be evaluated empirically rather than assumed.

\section{Conclusion}
We present a preference-conditioned design for evaluating fairness in personality-aware LLM recommendation. The central principle is simple but consequential: a changed recommendation is not automatically an unfair recommendation. By combining held-out preference utility, controlled demographic counterfactuals, measured Big Five profiles, matched negative controls, repeated generations, and multi-model/domain evaluation, the study is designed to separate beneficial personalization from unjustified identity-conditioned disparities. Final empirical conclusions will be inserted only after the frozen pipeline has generated and validated the corresponding evidence.

\bibliographystyle{splncs04}
\bibliography{references}

\end{document}
